\documentclass[11pt]{article}
\usepackage[utf8]{inputenc}
\usepackage[T1]{fontenc}
\usepackage[spanish]{babel}
\usepackage{amsmath, amssymb, bm}
\usepackage{geometry}
\geometry{margin=2.5cm}

\title{Ridge/LS 01: intuicion del operador regularizado}
\author{}
\date{}

\begin{document}
\maketitle

\section{Objetivo}

La expresion
\[
    (A^H A + \lambda I)^{-1}A^H
\]
aparece en el programa al estimar el canal a partir de las portadoras piloto.
Su funcion es construir un operador que toma mediciones conocidas del canal en
pilotos y devuelve una estimacion de los taps temporales del canal.

En el codigo, esta idea aparece en \texttt{core/ofdm\_ops.py}:
\[
    \texttt{weights}
    =
    (A^H A + \lambda I)^{-1}A^H
\]
y luego:
\[
    \hat{\bm h} = \texttt{weights}\,\bm y.
\]

\section{Que representan las letras}

El modelo simplificado es:
\[
    \bm y = A\bm h + \bm n.
\]

Aqui:

\begin{itemize}
    \item \(\bm y\) contiene las estimaciones de canal en las subportadoras piloto.
    \item \(A\) es una matriz que relaciona los taps temporales del canal con su respuesta en frecuencia en pilotos.
    \item \(\bm h\) contiene los taps temporales del canal.
    \item \(\bm n\) representa ruido y errores de medicion.
\end{itemize}

En el simulador, \(A\) se construye con exponenciales complejas:
\[
    A[m,\ell] = e^{-j2\pi k_m \ell/N}.
\]

Esto viene directamente de la DFT: un canal en tiempo produce una respuesta en
frecuencia.

\section{La pregunta que se intenta responder}

El receptor conoce \(\bm y\) y conoce \(A\). Quiere encontrar \(\bm h\).

Si no hubiera ruido y todo fuera perfectamente invertible, se podria pensar en:
\[
    \bm h = A^{-1}\bm y.
\]

Pero normalmente esto no se puede hacer directamente porque:

\begin{itemize}
    \item \(A\) no suele ser cuadrada.
    \item puede haber mas mediciones que incognitas, o menos.
    \item hay ruido.
    \item algunas columnas de \(A\) pueden parecerse demasiado entre si.
\end{itemize}

Entonces no buscamos una inversa exacta, sino una solucion que ajuste lo mejor
posible las mediciones.

\section{La idea de ajuste}

Queremos que:
\[
    A\hat{\bm h} \approx \bm y.
\]

Es decir, si usamos los taps estimados \(\hat{\bm h}\), deberiamos reconstruir
lo que observamos en los pilotos.

La diferencia:
\[
    \bm y - A\hat{\bm h}
\]
es el error de ajuste.

El metodo LS, o minimos cuadrados, busca que ese error sea pequeno.

\section{Por que aparece \(A^H\)}

La matriz \(A\) es compleja. Por eso no usamos solamente transpuesta \(A^T\),
sino transpuesta conjugada:
\[
    A^H.
\]

Si:
\[
    A =
    \begin{bmatrix}
        1 & j\\
        2 & 1-j
    \end{bmatrix},
\]
entonces:
\[
    A^H =
    \begin{bmatrix}
        1 & 2\\
        -j & 1+j
    \end{bmatrix}.
\]

En sistemas complejos, \(A^H\) es la herramienta natural para formar productos
internos, energias y ecuaciones normales.

\section{Que es \(A^H A\)}

La matriz:
\[
    A^H A
\]
se llama matriz de Gram.

Intuitivamente mide cuanto se parecen las columnas de \(A\) entre si.

\begin{itemize}
    \item En la diagonal aparecen las energias de las columnas.
    \item Fuera de la diagonal aparecen correlaciones entre columnas.
\end{itemize}

Si las columnas son muy diferentes e independientes, \(A^H A\) se comporta
bien. Si algunas columnas se parecen demasiado, \(A^H A\) puede volverse dificil
de invertir.

\section{Para que sirve \(\lambda I\)}

La parte:
\[
    \lambda I
\]
es una regularizacion.

La matriz identidad \(I\) tiene unos en la diagonal:
\[
    I =
    \begin{bmatrix}
        1 & 0 & 0\\
        0 & 1 & 0\\
        0 & 0 & 1
    \end{bmatrix}.
\]

Multiplicar por \(\lambda\) agrega \(\lambda\) a la diagonal de \(A^H A\):
\[
    A^H A + \lambda I.
\]

Esto hace que la matriz sea mas estable para invertir.

\section{Interpretacion sencilla}

La expresion:
\[
    (A^H A + \lambda I)^{-1}A^H
\]
es una especie de inversa corregida de \(A\).

No intenta invertir \(A\) de forma perfecta. Hace algo mas prudente:

\begin{itemize}
    \item ajusta las mediciones de pilotos;
    \item evita soluciones exageradas;
    \item reduce sensibilidad al ruido;
    \item permite estimar los taps del canal.
\end{itemize}

\section{Conexion con el programa}

En el simulador:

\[
    \hat{\bm h}
    =
    (A^H A + \lambda I)^{-1}A^H\bm y.
\]

Luego se calcula:
\[
    \hat{\bm H} = \mathrm{FFT}\{\hat{\bm h}\}.
\]

Finalmente se ecualiza:
\[
    \hat{X}[k] = \frac{Y[k]}{\hat{H}[k]}.
\]

Por tanto, el operador estudiado no ecualiza directamente. Primero estima el
canal en tiempo. Despues el programa obtiene la respuesta en frecuencia y recien
ahi ecualiza.

\end{document}
